\textbf{结论名称：} 向量模长公式

\textbf{适用条件：} 平面直角坐标系中，已知向量的坐标或起点、终点坐标。

\textbf{变量说明：} 设向量 $\vec{a}=(x,y)$，则其模长记为 $|\vec{a}|$，表示向量的长度（大小）。

\textbf{核心公式：} 
\[
\boxed{|\vec{a}| = \sqrt{x^{2} + y^{2}}}
\]
若已知起点 $A(x_1,y_1)$，终点 $B(x_2,y_2)$，则向量 $\overrightarrow{AB}$ 的模长为
\[
|\overrightarrow{AB}| = \sqrt{(x_2-x_1)^{2} + (y_2-y_1)^{2}} .
\]

\textbf{使用场景：} 求向量的长度、计算两点间距离、判断模的相等或倍数关系、以及三角形中位线等几何问题中的长度计算。

\textbf{简单例子：} 已知向量 $\vec{a}=(3,4)$，则其模长 $|\vec{a}| = \sqrt{3^{2}+4^{2}} = \sqrt{9+16}=5$。若点 $P(-1,2)$，$Q(2,-2)$，则 $|\overrightarrow{PQ}| = \sqrt{(2+1)^{2}+(-2-2)^{2}} = \sqrt{3^{2}+(-4)^{2}} = 5$。